\clearpage
\section{Certified filling under arbitrarily nested binders}
\label{sec:copying-nested-filling}
The abstraction case of addressed graph filling is now proved for the literal
Figure~16 translation. Combining it with the unary clause gives a certified
filling theorem for every original term context built from a hole, lambdas and
unary operations, with arbitrarily many enclosing binders. The argument term
is unrestricted: it may itself contain applications, binary operations and
boxes. Its hole annotation includes every bound and ambient variable available
there, and arbitrary atom environments are retained.

The theorem constructs a single address before inspecting the argument graph.
It proves that the address is sole, that its translated host is binding-free,
and that filling it with an original term translation reconstructs the original
translation of the filled term, up to recursive isomorphism and equality of
the depth indices. This last equality is not an equivalence between different
padding allowances. Parallel argument assembly and common-allowance padding
remain the structural cases needed for the unrestricted context theorem in
Appendix~B.1, item~3.

\subsection{Deep replacement through serial assembly}


Changing the body of a box in either component of a serial graph commutes
with serial assembly up to ordinary graph isomorphism. The box's bound,
free-variable and atom interfaces are fixed. Consequently the outer pattern
vertices and the two components' tagged vertex and edge carriers remain in
bijection: each tag and each local element is retained. The explicit maps
split on the component tag because the two dependent filling families are
not definitionally the same family of types.

For a left box, a left edge is either the selected occurrence, whose label
changes to the new body on both sides, or another occurrence, whose label is
unchanged. Every right label is unchanged. The right-box proof reverses those
roles. Both source and target lists retain their original component entries.
Mapping a component vertex through assembly gives the same result for an
internal vertex, an input port or an output port; the serial pattern identifies
the corresponding shared boundary in both graphs. Mapping these equalities
over the ordered lists proves incidence preservation. Equal-labelled boxes
at other occurrences are not changed.

\lean{Copying/SerialDeepPlug}{10}{61}{SerialDeepPlug.replaceBox_left}


\lean{Copying/SerialDeepPlug}{62}{113}{SerialDeepPlug.replaceBox_right}


An arbitrary addressed site can now be carried into either serial component.
For a shallow site, use the component edge's original slot tag and transport
the local interfaces along the already proved label equality. For a site
inside a box, tag the outer box edge and keep its body address unchanged.
This keeps the local depth and both local interfaces exactly as before.

To commute filling with the lifted site, split on the original site's first
constructor. A shallow site is handled by the preceding section's serial
replacement theorem; equality transport supplies the matching interfaces.
A site inside a box is handled by the new box-body interchange theorem,
using the fully filled inner body. No induction over the remaining address
is necessary at this step: serial assembly never changes that box's body
before filling. The resulting laws cover addresses at every finite nesting
depth on either side of serial composition.

\lean{Copying/SerialDeepPlug}{114}{122}{SerialDeepPlug.leftSite}


\lean{Copying/SerialDeepPlug}{123}{131}{SerialDeepPlug.rightSite}


\lean{Copying/SerialDeepPlug}{132}{141}{SerialDeepPlug.plugAt_left}


\lean{Copying/SerialDeepPlug}{142}{151}{SerialDeepPlug.plugAt_right}


Hole-free labels outside a chosen component edge remain hole-free in the
assembled graph, provided the other component is entirely hole-free. An
assembled edge either has the selected component tag, where inequality of
tagged edges implies inequality of local edges, or the other tag, where the
other component's hypothesis applies. This proves both left and right laws.

A sole-hole certificate then transports with the actual site. For a shallow
site, its marked label remains the same hole and the other-edge law certifies
all remaining edges. For a deep site, retain the inner sole-hole certificate
and use the other-edge law outside its enclosing box. These results certify
the addresses used by the filling laws themselves.

\lean{Copying/SerialDeepPlug}{152}{161}{SerialDeepPlug.left_other_free}


\lean{Copying/SerialDeepPlug}{162}{171}{SerialDeepPlug.right_other_free}


\lean{Copying/SerialDeepPlug}{172}{183}{SerialDeepPlug.leftSite_sole}


\lean{Copying/SerialDeepPlug}{184}{194}{SerialDeepPlug.rightSite_sole}


\subsection{Boxes and the full abstraction clause}
\termdef{copying-box-context-filling}


An equality cast on the host's outer interfaces transports its address
without changing the local hole interface or depth. Filling the transported
site equals casting the original filling's outer interfaces. Both definitions
and the equality proof follow by eliminating the two interface equalities.
This is the precise transport needed for the existing lambda-box convention,
which writes the bound variable first.

\lean{Copying/BoxContextPlug}{10}{13}{BoxContextPlug.Site.castHost}


\lean{Copying/BoxContextPlug}{14}{19}{BoxContextPlug.plugAt_castHost}


To enclose an address in a fresh box, select that box's singleton outer edge
and retain the original body address below it. Filling this new address is
isomorphic to boxing the filled body. Keep the empty outer internal carrier
and singleton edge carrier fixed. Every outer edge is the selected edge, so
its replacement label is exactly the box with the filled body. Sources and
targets agree because the bound and free interfaces are unchanged.

Recursive body isomorphisms also lift through a box. The same outer carrier
maps preserve incidence, while the recursive box-label relation is witnessed
by the given body isomorphism. This separates the exact box-construction
interchange from the later use of an isomorphism between body results.

\lean{Copying/BoxContextPlug}{20}{25}{BoxContextPlug.boxSite}


\lean{Copying/BoxContextPlug}{26}{43}{BoxContextPlug.box_filling}


\lean{Copying/BoxContextPlug}{44}{56}{BoxContextPlug.box_congr}


The sole-hole predicate is invariant under outer interface equality casts,
by equality elimination. Enclosing a sole site in a box preserves its
certificate: the inner certificate is unchanged and there is no other outer
edge, since the box's edge carrier is a singleton. Thus neither transport
nor boxing introduces an extra hole.

\lean{Copying/BoxContextPlug}{57}{61}{BoxContextPlug.SoleHole.castHost}


\lean{Copying/BoxContextPlug}{62}{71}{BoxContextPlug.boxSite_sole}


Outer interface casts preserve recursive isomorphism, again by equality
elimination. For the actual lambda box, first apply the existing cast from
$k+1$ variables to $1+k$ variables, then enclose the transported address in
the box. The box-filling isomorphism together with the host-cast filling
equality proves the law for this exact lambda-box constructor. The ordered
bound-variable convention is preserved; no alternative box interface is
introduced.

\lean{Copying/BoxContextPlug}{72}{77}{BoxContextPlug.castInterfaces_congr}


\lean{Copying/BoxContextPlug}{78}{81}{BoxContextPlug.lambdaBoxSite}


\lean{Copying/BoxContextPlug}{82}{89}{BoxContextPlug.lambdaBox_filling}


Follow the body address into its lambda box, through the right side of the
lambda-operation prefix, and through the left side of the prescribed
single-use distributor suffix. This defines the full abstraction address.

The filling law is a composition of three proved isomorphisms: commute filling
past the distributor suffix, commute it past the lambda prefix, and apply
the lambda-box filling law. Serial congruence lifts the inner isomorphisms
through the unchanged surrounding components. Abstraction congruence follows
from the same serial and box congruence laws, with the original interface
cast applied to the given body isomorphism. In particular, these proofs
retain both the lambda operation and all distributor edges.

\lean{Copying/BoxContextPlug}{90}{95}{BoxContextPlug.abstractionSite}


\lean{Copying/BoxContextPlug}{96}{103}{BoxContextPlug.abstraction_filling}


\lean{Copying/BoxContextPlug}{104}{108}{BoxContextPlug.abstraction_congr}


A unary constructor lifts any body address through its operation prefix and
distributor suffix, even when that address is below boxes. The two deep serial
filling laws prove the unary filling law directly. This strictly extends the
preceding section's shallow unary edge result.

For abstraction, sole-hole preservation follows in order from host-cast
preservation, singleton boxing, the hole-free lambda prefix, and the hole-free
distributor suffix. For a unary operation, use the corresponding hole-free
operation and suffix. Each constructor therefore supplies both the filling
isomorphism and the sole-hole proof for the same explicitly lifted address.

\lean{Copying/BoxContextPlug}{109}{114}{BoxContextPlug.unarySite}


\lean{Copying/BoxContextPlug}{115}{120}{BoxContextPlug.unary_filling}


\lean{Copying/BoxContextPlug}{121}{127}{BoxContextPlug.abstractionSite_sole}


\lean{Copying/BoxContextPlug}{128}{134}{BoxContextPlug.unarySite_sole}


\subsection{Induction over the original nested contexts}
\termdef{copying-scoped-isomorphism}


A scoped graph packages its finite allowance with its graph. To compare such
packages, define scoped isomorphism to mean an equality of their allowances
and a recursive graph isomorphism after transporting along that equality.
This merely handles dependent indices; it does not quotient different depths,
assert padding invariance, or alter any graph.

Reflexivity uses the identity recursive isomorphism. An isomorphism between
graphs at the same depth gives a scoped isomorphism immediately. For
transitivity, unpack the three depths, eliminate the two equalities, and
compose the recursive isomorphisms. At an already common depth, proof
irrelevance reduces any equality witness to reflexivity, so scoped
isomorphism yields an actual recursive isomorphism. Abstraction and unary
congruence follow by eliminating the depth equality and applying their
proved graph-constructor congruence laws.

\lean{Copying/NestedContextFilling}{11}{15}{NestedContextFilling.ScopedIso}


\lean{Copying/NestedContextFilling}{16}{17}{NestedContextFilling.refl}


\lean{Copying/NestedContextFilling}{18}{20}{NestedContextFilling.ofIso}


\lean{Copying/NestedContextFilling}{21}{33}{NestedContextFilling.trans}


\lean{Copying/NestedContextFilling}{34}{36}{NestedContextFilling.sameDepth}


\lean{Copying/NestedContextFilling}{37}{45}{NestedContextFilling.abstraction}


\lean{Copying/NestedContextFilling}{46}{59}{NestedContextFilling.unary}


The allowed-context predicate is an inductive predicate on the original
annotated context grammar: a hole is allowed, an enclosing lambda preserves
the predicate, and any enclosing unary operation preserves it. It introduces
no new term syntax and imposes no restriction on the eventual argument term.
It intentionally does not yet cover application or binary-operation context
constructors.

The literal hole corolla has a shallow address at its annotated term and
environment interface at every finite allowance. Its sole-hole proof uses
the singleton edge carrier. Filling that address is recursively isomorphic
to the supplied graph by the generic corolla left-unit isomorphism.

\lean{Copying/NestedContextFilling}{60}{67}{NestedContextFilling.UnaryLambda}


\lean{Copying/NestedContextFilling}{68}{71}{NestedContextFilling.holeSite}


\lean{Copying/NestedContextFilling}{72}{76}{NestedContextFilling.holeSite_sole}


\lean{Copying/NestedContextFilling}{77}{85}{NestedContextFilling.hole_filling}


Induct on the allowed-context derivation. The induction statement constructs
one actual address in the literal context graph, at the original hole's
allowance, proves that it is sole, and proves the filling correspondence
for \emph{every} graph at that interface. The address therefore does not
depend on the eventual argument.

The hole case is the corolla law. For a lambda, lift the induction address
through the abstraction constructor. Its sole-hole certificate follows from
the abstraction certificate. Compose the abstraction filling isomorphism
with abstraction congruence applied to the induction hypothesis. For a unary
operation, use the analogous unary address, certificate, filling law and
congruence. The resulting target is exactly the existing constructor
interpretation at the supplied argument graph.

Finally choose the hole allowance to be the original argument term's
translation depth and supply its original translated graph. The existing
constructor-interpretation theorem identifies the target with translation
of the syntactically filled term. The independent binding-free theorem
certifies the literal host context graph. This proves the final conjunction:
an actual address, its sole-hole proof, host binding-freedom, and scoped
isomorphism between actual filling and the unchanged filled-term translation.
The proof does not appeal to execution or observational equivalence.

\lean{Copying/NestedContextFilling}{86}{106}{NestedContextFilling.unaryLambda_filling}


\lean{Copying/NestedContextFilling}{107}{120}{NestedContextFilling.unaryLambda_term_filling}


The regression context is
$\lambda.\,\mathtt{ref}(\lambda.\,![-])$.
Its hole has two more variable slots than its outer environment and the
same atom slots. The allowed-context derivation follows the two lambdas and
two unary operations in order. A reflexivity theorem checks the exact filled
term, with neither binder nor operation moved or erased.

For every ambient environment size, hole name and argument term at that
annotated scope, the final regression obtains a sole-hole, binding-free host
and the actual graph-filling correspondence with the direct translation of
$\lambda.\,\mathtt{ref}(\lambda.\,!t)$.
This exercises address transport below two boxes with operations between
and below them. It is a structural regression; it makes no claim about
executing reference creation or dereference on the resulting term.

\lean{Copying/NestedFillingExamples}{7}{9}{NestedFillingExamples.twoBinderContext}


\lean{Copying/NestedFillingExamples}{10}{13}{NestedFillingExamples.twoBinderContext_allowed}


\lean{Copying/NestedFillingExamples}{14}{18}{NestedFillingExamples.twoBinderContext_fill}


\lean{Copying/NestedFillingExamples}{19}{28}{NestedFillingExamples.twoBinderContext_graph_filling}


All new declarations pass the full build, and the three axiom audits report
only standard Lean axioms. This closes the standalone abstraction and nested
lambda/unary filling cases. Item~3 still needs parallel argument assembly,
compatible padding, the unrestricted context induction, and term-context
congruence. Execution's isomorphism and padding laws must be discharged by
the copying machine. Proposition~12.2 and its translation remain unchanged.
